% ----------------------
\subsection{Section 113 (Between circa 16 and circa 29 March 1838)}
% ----------------------
	\label{DC113}
	
	% -----
	\subsubsection{Historical Background}
	% -----
		\label{DC113-HB}
		
		% --
		\paragraph{JSP}
		% --
		
			%
			\subparagraph{Verses 1--6}
			%
			
				``Sometime in mid- or late March 1838, after JS arrived in Far West, Missouri, he apparently answered a series of questions regarding the prophecies in Isaiah chapters 11 and 52. The questions and answers were inscribed in JS's `Scriptory Book' by George Robinson. No authorship was attributed to the questions and answers, although some of the answers begin with `thus saith the Lord.' JS's authorship is implied, in that some of the answers are couched in the language of revelation, similar to other revelations transcribed in the Scriptory Book. Further, most of the documents transcribed in the book are explicitly JS documents.

				``When the questions and answers regarding the prophecies of Isaiah were copied into the multivolume manuscript history JS began in 1838, they were grouped under the heading `Questions on Scripture.' Following this format, the questions and answers were presented as a unified set when the history was later printed in the church's newspaper. However, the five question-and-answer pairs were inscribed in the Scriptory Book under two headings: the first three pairs appear under the heading `Quest. on Scripture,' while the remaining two appear under the heading `Questions by Elias Higby.' These headings are comparable to the other headings in the Scriptory Book that were used to demarcate different document transcripts. The first three questions, posed by an unidentified interlocutor, regard Isaiah chapter 11, while the other two questions, posed by Elias Higbee, concern Isaiah chapter 52. Additionally, whereas the first three questions were answered with `thus saith the Lord' responses, the other questions were not. These differences further suggest what seems to be indicated by the distinct headings: that the two sets of question-and-answer pairs were considered separate, albeit obviously related, texts.

				``The dating of these two texts is uncertain, but their location in the Scriptory Book suggests an approximate period of between 16 and 29 March 1838. The questions and answers follow a motto for the church, which JS composed at the earliest on 16 March, and are followed by several document transcripts that are arranged in roughly chronological order. These documents begin with two that JS produced in September 1837 in Kirtland, before he migrated to Far West. This pair of Kirtland documents raises the possibility that some or all of the questions and answers also date back to the Kirtland period. However, the similarities in the format and the subject of the two sets of questions suggest that all of the questions were posed and answered in the same setting. Further, the fact that some of the questions were posed by Elias Higbee---who had been living in Caldwell County, Missouri, since 1836 and had been serving on the Zion high council in Far West---indicates that the questions and answers date to the period after JS arrived in Far West on 14 March. Following the pair of Kirtland documents is JS's letter to the Kirtland presidency on 29 March 1838, strongly suggesting that the questions and answers were composed between 16 March and 29 March 1838. As with the other documents inscribed in the Scriptory Book, Robinson probably transcribed the questions and answers from an earlier manuscript. If the questions and answers were composed following Robinson's arrival in Far West on 28 March, he could have dated the documents, which suggests they were composed earlier. While JS could have met with Higbee anytime between 16 and 29 March, the only known time they were together was during a meeting of the high council and bishopric on 24 March 1838, suggesting the possibility that the discussion took place after the meeting adjourned that evening.

				``Isaiah chapter 11, the subject of the first set of questions, figures prominently in JS's revelations and writings. Within a few months of answering these questions, JS recounted in his manuscript history that when visited by the angel Moroni in 1823, Moroni `quoted the Eleventh Chapter of Isaiah saying that it was about to be fulfilled.' The Book of Mormon draws heavily on Isaiah's prophetic worldview and encourages readers to `search the prophecies of Isaiah.' Several times, the volume singles out the `great' writings of Isaiah from the writings of other Israelite prophets, and the volume quotes extensively from the book of Isaiah, including entire chapters. For example, the Book of Mormon quotes Isaiah chapter 11 in its entirety and then quotes verses 4--9 again later on in the volume. Furthermore, JS's subsequent revelations and writings repeat or allude to several verses in this chapter, applying them to the gathering of Israel in the last days. Shortly before JS left Kirtland, a member of the First Quorum of the Seventy delivered a sermon there `on the gathering of the house of Israel,' using Isaiah 11 as his text.

				``Isaiah 11:1 contains the prophecy that a `rod' or a `branch' (a new shoot) would grow out of the `stem' (stump) or `roots' of Jesse, the father of the Israelite king David. The next four verses have traditionally been interpreted as further prophecy regarding what this David-like messianic figure would do for the children of Israel. However, the first two questions apparently assume that these verses describe the stem, or old stump, of Jesse rather than the rod or new branch growing out of the stem. The first question regards the identity of the stem of Jesse. Whereas the Christian exegetical tradition generally interprets the stem of Jesse as the Davidic dynasty, JS interpreted it as Jesus Christ. In addition, whereas most Christians interpret the root of Jesse to be Christ, JS interpreted the root as a latter-day figure. As evidenced in JS's earlier revelations regarding the Bible, including his revision, or `new translation,' of the Bible, he considered himself a prophet similar to those in the Old Testament, with full authority to receive new revelation to interpret and clarify the writings of his predecessors. The answers to the second and third questions, regarding the rod coming out of the stem of Jesse and the root of Jesse, suggest that JS or a similar latter-day figure would fulfill these prophecies. These answers came within the context of JS experiencing the fallout of apostasy in Kirtland and then engaging in a reassertion of his prophetic authority in the land of Zion.

				``Whereas the second set of questions was posed by Higbee, it is uncertain who posed the first set of questions. Although presumably not Higbee, it could have been another person present at the 24 March meeting, or it could have been JS, petitioning the Lord for revealed answers to his questions. It is also unknown who originally wrote down the questions and answers, though JS's revelations were usually written by someone else acting in the role of scribe. Robinson probably had not yet arrived in Far West. Higbee, who served as a judge in Caldwell County and was soon called as a church historian, may have had some clerical duties. Although Robinson began the Scriptory Book in late March, with an account of JS's arrival in Far West and a copy of the motto, he apparently did not add anything further to the book until mid-April. Therefore, Robinson likely copied the questions and answers from a loose manuscript into the Scriptory Book, probably sometime in mid- or late April.''
			
			%
			\subparagraph{Verses 7--10}
			%
			
				``Sometime after JS arrived in Far West, Missouri, he apparently answered two sets of questions, one regarding Isaiah chapter 11 and the other, labeled `Questions by Elias Higby [Higbee],' relating to Isaiah chapter 52. The dating of these two sets of questions and answers is uncertain, but they were probably produced sometime between 16 and 29 March 1838. While JS could have met with Higbee, who was a member of the Zion high council, anytime during this period, the one time they were known to be together was during a meeting of the high council and bishopric on 24 March 1838, suggesting that Higbee may have posed his questions to JS after the meeting adjourned that evening. The answers to the first set of questions, regarding Isaiah 11, begin with `thus saith the Lord.' The revealed answers may have provoked Higbee to ask the questions about Isaiah 52:1--2, though the answers given to Higbee's questions do not include the same revelatory language. These verses regard the redemption of Zion and introduce Isaiah's suffering servant oracles, including the material in chapter 53, which is one of traditional Christianity's most important texts for Christianizing the Old Testament. Isaiah 52:1--2 is quoted twice in the Book of Mormon, and other passages in chapter 52 also appear in the book. JS revelations also repeat or allude to verses in this chapter, using them to explain the gathering and reestablishment of the house of Israel in the latter days.
				
				``Higbee posed two questions. The answers, presumably from JS, stated that those called of God in the last days would be given priesthood authority to gather scattered Israel home to Zion. These answers may have been considered as having special significance at this time because of JS's move to Missouri, the Saints' `Land of Zion.' Probably sometime in mid- or late April, George W. Robinson inscribed the questions and answers in the `Scriptory Book,' probably from an earlier manuscript.''
					
	% -----
	\subsubsection{Versions}
	% -----
		\label{DC113-V}
		
		See the \href{https://drive.google.com/file/d/1goAvNz8jS4R8slYfMG_db55Ty2z_vmgK/view?usp=sharing}{sections comparison} document.
	
		\begin{compactdesc}
			\item[JSP] \phantom{.}
				\begin{compactdesc}
					\item[Verses 1--6] \hyperref[EMS-1838-JS-CIR-16-CIR-29-MAR-1838-A]{Between circa 16 and circa 29 March 1838}
					\item[Verses 7--10] \hyperref[EMS-1838-JS-CIR-16-CIR-29-MAR-1838-B]{Between circa 16 and circa 29 March 1838}
				\end{compactdesc}
			\item[BC] ---
			\item[DC 1835] ---
			\item[DC 1844] ---
			\item[TS] ---
			\item[DN] 3:8 (5 March 1853), 29
			\item[MS] 16:8 (25 February 1854), 117--118
		\end{compactdesc}

	% -----
	\subsubsection{Section 113:1} % *DC113-1 
	% -----
		\label{DC113-1}

		WHO is the Stem of Jesse spoken of in the \hyperref[ISAIAH-11-1]{1st}, \hyperref[ISAIAH-11-2]{2d}, \hyperref[ISAIAH-11-3]{3d}, \hyperref[ISAIAH-11-4]{4th}, and \hyperref[ISAIAH-11-5]{5th} verses of the 11th chapter%
			\footnote{%
				% \label{}%
				See \hyperref[ISAIAH-11]{Isaiah 11}.
				}
		of Isaiah?

		% \begin{framed}
		% \scriptsize
			% \begin{compactdesc}
				% \item[JSP] 
				% % \item[BC] 
				% % \item[]
			% \end{compactdesc}
		% \end{framed}

	% -----
	\subsubsection{Section 113:2} % *DC113-2 
	% -----
		\label{DC113-2}

		Verily thus saith the Lord: It is Christ.

		% \begin{framed}
		% \scriptsize
			% \begin{compactdesc}
				% \item[JSP] 
				% % \item[BC] 
				% % \item[]
			% \end{compactdesc}
		% \end{framed}

	% -----
	\subsubsection{Section 113:3} % *DC113-3 
	% -----
		\label{DC113-3}

		What is the rod spoken of in the first verse of the 11th chapter of Isaiah, that should come of the Stem of Jesse?

		% \begin{framed}
		% \scriptsize
			% \begin{compactdesc}
				% \item[JSP] 
				% % \item[BC] 
				% % \item[]
			% \end{compactdesc}
		% \end{framed}

	% -----
	\subsubsection{Section 113:4} % *DC113-4 
	% -----
		\label{DC113-4}

		Behold, thus saith the Lord: It is a servant in the hands of Christ, who is partly a descendant of Jesse as well as of Ephraim, or of the house of Joseph, on whom there is laid much power.

		% \begin{framed}
		% \scriptsize
			% \begin{compactdesc}
				% \item[JSP] 
				% % \item[BC] 
				% % \item[]
			% \end{compactdesc}
		% \end{framed}

	% -----
	\subsubsection{Section 113:5} % *DC113-5 
	% -----
		\label{DC113-5}

		What is the root of Jesse spoken of in the 10th verse of the 11th chapter?

		% \begin{framed}
		% \scriptsize
			% \begin{compactdesc}
				% \item[JSP] 
				% % \item[BC] 
				% % \item[]
			% \end{compactdesc}
		% \end{framed}

	% -----
	\subsubsection{Section 113:6} % *DC113-6 
	% -----
		\label{DC113-6}

		Behold, thus saith the Lord, it is a descendant of Jesse, as well as of Joseph, unto whom rightly belongs the priesthood, and the keys of the kingdom, for an ensign, and for the gathering of my people in the last days.

		% \begin{framed}
		% \scriptsize
			% \begin{compactdesc}
				% \item[JSP] 
				% % \item[BC] 
				% % \item[]
			% \end{compactdesc}
		% \end{framed}

	% -----
	\subsubsection{Section 113:7} % *DC113-7 
	% -----
		\label{DC113-7}

		Questions by Elias Higbee: What is meant by the command in Isaiah, 52d chapter, \hyperref[ISAIAH-52-1]{1st} verse, which saith: Put on thy strength, O Zion---and what people had Isaiah reference to?

		% \begin{framed}
		% \scriptsize
			% \begin{compactdesc}
				% \item[JSP] 
				% % \item[BC] 
				% % \item[]
			% \end{compactdesc}
		% \end{framed}

	% -----
	\subsubsection{Section 113:8} % *DC113-8 
	% -----
		\label{DC113-8}

		He had reference to those whom God should call in the last days, who should hold the power of priesthood to bring again Zion, and the redemption of Israel; and to put on her strength is to put on the authority of the priesthood, which she, Zion, has a right to by lineage; also to return to that power which she had lost.

		% \begin{framed}
		% \scriptsize
			% \begin{compactdesc}
				% \item[JSP] 
				% % \item[BC] 
				% % \item[]
			% \end{compactdesc}
		% \end{framed}

	% -----
	\subsubsection{Section 113:9} % *DC113-9 
	% -----
		\label{DC113-9}

		What are we to understand by Zion loosing herself from the bands of her neck; 2d verse?

		% \begin{framed}
		% \scriptsize
			% \begin{compactdesc}
				% \item[JSP] 
				% % \item[BC] 
				% % \item[]
			% \end{compactdesc}
		% \end{framed}

	% -----
	\subsubsection{Section 113:10} % *DC113-10 
	% -----
		\label{DC113-10}

		We are to understand that the scattered remnants are exhorted to return to the Lord from whence they have fallen; which if they do, the promise of the Lord is that he will speak to them, or give them revelation.  See the 6th, 7th, and 8th verses.  The bands of her neck are the curses of God upon her, or the remnants of Israel in their scattered condition among the Gentiles.

		% \begin{framed}
		% \scriptsize
			% \begin{compactdesc}
				% \item[JSP] 
				% % \item[BC] 
				% % \item[]
			% \end{compactdesc}
		% \end{framed}